%% TCC -- Apresentação de Defesa
%% Previsão de Insuficiência Cardíaca com Redes Neurais Artificiais
%% Lucas | UFVJM | Orientador: Prof. Dr. Leonardo Lana de Carvalho
%% Compilar com: pdflatex apresentacao.tex

\documentclass[aspectratio=169,11pt,brazil]{beamer}

\usepackage[utf8]{inputenc}
\usepackage[T1]{fontenc}
\usepackage[brazil]{babel}
\usepackage{mathptmx}
\usepackage{graphicx}
\usepackage{booktabs}
\usepackage{xcolor}
\usepackage{pgfplots}
\pgfplotsset{compat=1.18}
\usepackage{tikz}
\usetikzlibrary{positioning}
\usepackage{amsmath}
\usepackage{array}
\usepackage{colortbl}

% -----------------------------------------------------------------------
% Identidade visual — azul claro UFVJM
% -----------------------------------------------------------------------
\definecolor{ufvjmblue}{RGB}{0,100,180}
\definecolor{ufvjmbluedark}{RGB}{0,70,130}
\definecolor{ufvjmbluelight}{RGB}{200,225,245}
\definecolor{ufvjmgray}{RGB}{90,90,90}

% Cores dos modelos nos gráficos
\definecolor{cAdalineNP}{RGB}{0,100,180}
\definecolor{cAdalineSK}{RGB}{80,170,230}
\definecolor{cMlpNP}{RGB}{210,70,30}
\definecolor{cMlpSK}{RGB}{255,165,30}
\definecolor{cBin}{RGB}{0,100,180}
\definecolor{cCont}{RGB}{255,165,30}

\usetheme{Madrid}
\usecolortheme[named=ufvjmblue]{structure}
\setbeamercolor{palette primary}{bg=ufvjmblue,fg=white}
\setbeamercolor{palette secondary}{bg=ufvjmblue!80,fg=white}
\setbeamercolor{palette tertiary}{bg=ufvjmblue!60,fg=white}
\setbeamercolor{title}{bg=ufvjmblue,fg=white}
\setbeamercolor{subtitle}{fg=white}
\setbeamercolor{frametitle}{bg=ufvjmblue,fg=white}
\setbeamercolor{block title}{bg=ufvjmblue,fg=white}
\setbeamercolor{block body}{bg=ufvjmbluelight,fg=black}
\setbeamercolor{alerted text}{fg=ufvjmbluedark}
\setbeamertemplate{navigation symbols}{}
\setbeamertemplate{caption}{\insertcaption}
\setbeamerfont{frametitle}{series=\bfseries}

% -----------------------------------------------------------------------
% Metadados
% -----------------------------------------------------------------------
\title[Uso de Redes Neurais para Diagnóstico Cardíaco]{Uso de Redes Neurais Artificiais para\\Diagnóstico de Doenças Cardíacas}
\author[Lucas dos Santos Martins]{Lucas dos Santos Martins}
\institute[UFVJM]{%
  Orientador: Prof.\ Dr.\ Leonardo Lana de Carvalho\\[0.2em]
  Banca: Prof.\ André Luiz Covre \quad Prof.\ Áthila Rocha Trindade\\[0.4em]
  Universidade Federal dos Vales do Jequitinhonha e Mucuri}
\date{2026}

\begin{document}

% =======================================================================
% SLIDE 1 — CAPA
% =======================================================================
\begin{frame}[plain]
  \titlepage
\end{frame}

% =======================================================================
% SLIDE 2 — CONTEXTUALIZAÇÃO
% =======================================================================
\section{Contextualização}

\begin{frame}{Contextualização}
  \begin{columns}[T]
    \begin{column}{0.55\textwidth}
      \begin{itemize}
        \item Doenças cardiovasculares: \textbf{principal causa de morte no mundo}
        \item Diagnóstico precoce é fundamental, mas \textbf{nem sempre disponível a tempo}
      \end{itemize}
    \end{column}
    \begin{column}{0.45\textwidth}
      \begin{block}{Papel das Redes Neurais}
        Auxiliar o médico, não substituir o julgamento clínico, oferecendo de forma rápida e a qualquer momento uma \textbf{análise baseada nos dados do paciente}.
      \end{block}
      \vspace{0.4em}
      \begin{alertblock}{Impacto prático}
        Com base no resultado, o médico pode \textbf{priorizar investigações} e ficar mais atento a casos de risco.
      \end{alertblock}
    \end{column}
  \end{columns}
\end{frame}

% =======================================================================
% SLIDE 3 — OBJETIVO
% =======================================================================
\section{Objetivo}

\begin{frame}{Objetivo}
  \begin{block}{Objetivo Geral}
    Desenvolver e comparar modelos de \textbf{redes neurais artificiais} para a previsão de doenças cardíacas em dois contextos clínicos distintos.
  \end{block}
  \vspace{0.4em}
  \begin{columns}[T]
    \begin{column}{0.48\textwidth}
      \textbf{Modelos implementados:}
      \begin{itemize}
        \item \textbf{ADALINE Logística}: neurônio único
        \item \textbf{Perceptron Multicamadas (MLP)}: rede multicamadas
      \end{itemize}
      \smallskip
      Cada modelo em \textbf{duas versões}:\\
      \quad implementação manual\\
      \quad \textit{vs} biblioteca (Scikit-learn)
    \end{column}
    \begin{column}{0.52\textwidth}
      \textbf{Questões respondidas:}
      \begin{itemize}
        \item A complexidade extra do Multicamadas traz \textbf{ganho real}?
        \item As implementações manuais \textbf{reproduzem} os resultados da biblioteca?
        \item Qual modelo é mais \textbf{clinicamente seguro} (menos falsos negativos)?
        \item A \textbf{binarização prévia} dos dados traz ganhos para os modelos ou não?
      \end{itemize}
    \end{column}
  \end{columns}
\end{frame}

% =======================================================================
% SLIDE 4 — BANCOS DE DADOS
% =======================================================================
\section{Dados e Métodos}

\begin{frame}{Bancos de Dados}
  \begin{columns}[T]
    \begin{column}{0.50\textwidth}
      \begin{block}{Banco 1 — Heart Disease Dataset}
        \begin{itemize}
          \item \textbf{918 pacientes} — 11 variáveis clínicas\\
                (508 doentes | 410 saudáveis)
          \item Idade, pressão arterial, colesterol, tipo de dor torácica, eletrocardiograma (ECG)\ldots
          \item \textbf{Alvo:} presença de doença cardíaca (0/1)
          \item Divisão \textbf{70/30}:\\
                642 treino | 276 teste
        \end{itemize}
      \end{block}
    \end{column}
    \begin{column}{0.50\textwidth}
      \begin{block}{Banco 2 — Heart Failure Clinical Records}
        \begin{itemize}
          \item \textbf{299 pacientes} — 12 variáveis\\
                (96 óbitos | 203 sobreviventes)
          \item Idade, tempo de acompanhamento, anemia, diabetes, hipertensão\ldots
          \item \textbf{Alvo:} óbito durante acompanhamento (0/1)
          \item Divisão \textbf{70/30}:\\
                209 treino | 90 teste
        \end{itemize}
      \end{block}
    \end{column}
  \end{columns}
  \vfill
  \centering\small
  Contextos distintos: \textbf{diagnóstico de doença cardíaca} \textit{vs} \textbf{prognóstico de insuficiência cardíaca}
\end{frame}

% =======================================================================
% SLIDE 5 — PRÉ-PROCESSAMENTO
% =======================================================================
\begin{frame}{Pré-Processamento}
  \begin{columns}[T]
    \begin{column}{0.50\textwidth}
      \textbf{Transformações aplicadas (Banco 1):}
      \begin{itemize}
        \item \textbf{Separação em colunas binárias:} tipo de dor torácica\\
              ATA, NAP, ASY, TA $\to$ 4 colunas binárias
        \item \textbf{Mapeamento binário:}\\
              Sexo (M=1 / F=0)\\
              Angina (Y=1 / N=0)
        \item \textbf{Binarização por faixas clínicas:}\\
              Idade, colesterol, pressão arterial\ldots
      \end{itemize}
    \end{column}
    \begin{column}{0.50\textwidth}
      \textbf{Divisão treino/teste:}
      \begin{itemize}
        \item Proporção \textbf{70/30}
        \item Preserva proporção de casos positivos em cada partição
      \end{itemize}
    \end{column}
  \end{columns}
\end{frame}

% =======================================================================
% SLIDE 6 — ADALINE TEORIA
% =======================================================================
\section{Modelos}

\begin{frame}{ADALINE Logístico — Modelo e Justificativa}
  \begin{columns}[T]
    \begin{column}{0.50\textwidth}
      \textbf{O que é o ADALINE Logístico?}
      \begin{itemize}
        \item Modelo de \textbf{neurônio único}
        \item Aplica a função \textbf{sigmoide} na saída para classificação binária, retornando uma \textbf{probabilidade} entre 0 e 1
      \end{itemize}
    \end{column}
    \begin{column}{0.50\textwidth}
      \begin{block}{Por que usar o ADALINE Logístico?}
        \begin{itemize}
          \item \textbf{Referência simples:} se já performa bem, questiona-se o ganho real do Multicamadas
          \item Computacionalmente leve
        \end{itemize}
      \end{block}

    \end{column}
  \end{columns}
\end{frame}

% =======================================================================
% SLIDE 7 — MLP TEORIA
% =======================================================================
\begin{frame}{Perceptron Multicamadas (MLP) — Modelo e Justificativa}
  \begin{columns}[T]
    \begin{column}{0.50\textwidth}
      \textbf{O que é o Multicamadas?}
      \begin{itemize}
        \item Rede com \textbf{pelo menos uma camada oculta}
        \item Aprende relações \textbf{não-lineares} entre variáveis
        \item Nossa implementação:\\
              \textbf{10 neurônios ocultos}\\
              + 1 neurônio de saída (sigmoide)
      \end{itemize}
    \end{column}
    \begin{column}{0.50\textwidth}
      \begin{block}{Por que usar o Multicamadas?}
        \begin{itemize}
          \item Responde à pergunta central: a \textbf{complexidade extra traz ganho real} em dados clínicos?
          \item Representa o próximo nível acima do ADALINE Logístico
        \end{itemize}
      \end{block}
    \end{column}
  \end{columns}
\end{frame}

% =======================================================================
% SLIDE 8 — CONFIGURAÇÕES DE EXPERIMENTO
% =======================================================================
\section{Experimentos}

\begin{frame}{Configurações de Experimento}
  \begin{columns}[T]
    \begin{column}{0.55\textwidth}
      \textbf{Épocas:}
      \begin{itemize}
        \item Quantas vezes o modelo percorre os dados de treino
        \item Poucas épocas → modelo \textbf{não aprende o suficiente}
        \item Épocas demais → modelo \textbf{memoriza} os dados e perde capacidade de generalizar (\textit{overfitting})
      \end{itemize}
      \vspace{0.4em}
      \textbf{Taxa de aprendizado:}
      \begin{itemize}
        \item Controla o \textbf{tamanho do passo} de cada ajuste
        \item Taxa alta → aprende rápido, mas de forma \textbf{grosseira}
        \item Taxa baixa → ajustes mais \textbf{finos}, mas exige mais tempo
      \end{itemize}
    \end{column}
    \begin{column}{0.45\textwidth}
      \begin{block}{Três configurações por modelo}
        \vspace{0.1em}
        \centering
        \renewcommand{\arraystretch}{1.3}
        \begin{tabular}{ccc}
          \toprule
          \textbf{Config.} & \textbf{Épocas} & \textbf{Taxa} \\
          \midrule
          C1 & 500     & 0,01    \\
          C2 & 1.000   & 0,001   \\
          C3 & 3.000   & 0,0001  \\
          \bottomrule
        \end{tabular}
      \end{block}
      \vspace{0.25em}
      \small
      \begin{itemize}\setlength{\itemsep}{2pt}
        \item \textbf{C1} — rápido, ajustes grosseiros
        \item \textbf{C2} — equilíbrio velocidade/precisão
        \item \textbf{C3} — ajustes finos, risco de \textit{overfitting}
      \end{itemize}
    \end{column}
  \end{columns}
\end{frame}

% =======================================================================
% SLIDE 9 — RESULTADOS BANCO 1 (texto)
% =======================================================================
\section{Resultados}

\begin{frame}{Resultados — Banco 1: Heart Disease}
  \centering\small
  \textit{918 pacientes: 508 doentes | 410 saudáveis}
  \vspace{0.3em}

  \renewcommand{\arraystretch}{1.35}
  \begin{tabular}{lllll}
    \toprule
    \textbf{Modelo} & \textbf{C1} & \textbf{C2} & \textbf{C3} & \textbf{FN mín.} \\
    \midrule
    Adaline Sem Biblioteca & 85,87\% / 28\,FN & 86,23\% / 21\,FN & \textbf{86,96\% / 20\,FN} & 20 (C3) \\
    Adaline Com Biblioteca & 85,87\% / 21\,FN & 86,23\% / 21\,FN & \textbf{86,23\% / 21\,FN} & 21 (Todos) \\
    Multicamadas Sem Biblioteca     & 82,97\% / 27\,FN & 86,23\% / 16\,FN & \cellcolor{ufvjmbluelight}\textbf{87,68\% / 15\,FN} & \textbf{15 (C3)} \\
    Multicamadas Com Biblioteca     & 87,32\% / 18\,FN & \cellcolor{ufvjmbluelight}\textbf{88,04\% / 20\,FN} & 82,61\% / 35\,FN & 18 (C1) \\
    \bottomrule
  \end{tabular}

  \vfill
  \begin{columns}[T]
    \begin{column}{0.50\textwidth}
      \begin{block}{Maior acurácia}
        Multicamadas Com Biblioteca C2: \textbf{88,04\%}
      \end{block}
    \end{column}
    \begin{column}{0.50\textwidth}
      \begin{alertblock}{Melhor clinicamente}
        Multicamadas Sem Biblioteca C3: \textbf{15 FN}
      \end{alertblock}
    \end{column}
  \end{columns}
  \vspace{0.2em}
  \centering{\small\textbf{Melhor modelo geral:} Multicamadas Sem Biblioteca C3 — melhor equilíbrio entre acurácia (87,68\%) e FN mínimo (15)}
\end{frame}

% =======================================================================
% SLIDE 10 — RESULTADOS BANCO 1 (gráficos)
% =======================================================================
\begin{frame}{Banco 1 — Visualização dos Resultados}
  \begin{columns}[c]
    % ---- gráfico acurácia ----
    \begin{column}{0.50\textwidth}
      \centering
      {\small\textbf{Acurácia por Configuração (\%)}}
      \vspace{0.15em}

      \begin{tikzpicture}
        \begin{axis}[
          ybar,
          width=1.05\columnwidth,
          height=5.6cm,
          symbolic x coords={C1,C2,C3},
          xtick=data,
          xticklabel style={font=\scriptsize},
          ymin=79, ymax=92,
          ytick={80,82,84,86,88,90,92},
          yticklabel={\pgfmathprintnumber{\tick}},
          ylabel={\%},
          ylabel style={font=\scriptsize,rotate=-90,at={(ticklabel cs:0.5)},anchor=south},
          yticklabel style={font=\scriptsize},
          bar width=5.5pt,
          legend style={at={(0.5,-0.22)},anchor=north,legend columns=2,font=\tiny,
                        draw=none,fill=none,column sep=2pt},
          legend cell align=left,
          every axis plot/.append style={draw=none},
          axis line style={draw=ufvjmgray,thin},
          tick style={draw=ufvjmgray,thin},
          ymajorgrids=true,
          grid style={dotted,gray!40},
          clip=false,
          nodes near coords={\pgfmathprintnumber[fixed,precision=1]{\pgfplotspointmeta}},
          nodes near coords style={font=\tiny,inner sep=0pt,rotate=90,anchor=west,xshift=2pt},
        ]
          \addplot[fill=cAdalineNP] coordinates {(C1,85.87) (C2,86.23) (C3,86.96)};
          \addplot[fill=cAdalineSK] coordinates {(C1,85.87) (C2,86.23) (C3,86.23)};
          \addplot[fill=cMlpNP]     coordinates {(C1,82.97) (C2,86.23) (C3,87.68)};
          \addplot[fill=cMlpSK]     coordinates {(C1,87.32) (C2,88.04) (C3,82.61)};
          \legend{ADALINE SB, ADALINE CB, Multicamadas SB, Multicamadas CB}
        \end{axis}
      \end{tikzpicture}

      \vspace{0.2em}
      {\tiny $\uparrow$ maior acurácia = melhor desempenho geral}
    \end{column}
    % ---- gráfico FN ----
    \begin{column}{0.50\textwidth}
      \centering
      {\small\textbf{Falsos Negativos por Configuração}}
      \vspace{0.15em}

      \begin{tikzpicture}
        \begin{axis}[
          ybar,
          width=1.05\columnwidth,
          height=5.6cm,
          symbolic x coords={C1,C2,C3},
          xtick=data,
          xticklabel style={font=\scriptsize},
          ymin=0, ymax=43,
          ytick={0,10,20,30,40},
          ylabel={FN},
          ylabel style={font=\scriptsize,rotate=-90,at={(ticklabel cs:0.5)},anchor=south},
          yticklabel style={font=\scriptsize},
          bar width=5.5pt,
          legend style={at={(0.5,-0.22)},anchor=north,legend columns=2,font=\tiny,
                        draw=none,fill=none,column sep=2pt},
          legend cell align=left,
          every axis plot/.append style={draw=none},
          axis line style={draw=ufvjmgray,thin},
          tick style={draw=ufvjmgray,thin},
          ymajorgrids=true,
          grid style={dotted,gray!40},
          nodes near coords,
          nodes near coords style={font=\tiny,inner sep=1pt},
          nodes near coords align=above,
        ]
          \addplot[fill=cAdalineNP] coordinates {(C1,28) (C2,21) (C3,20)};
          \addplot[fill=cAdalineSK] coordinates {(C1,21) (C2,21) (C3,21)};
          \addplot[fill=cMlpNP]     coordinates {(C1,27) (C2,16) (C3,15)};
          \addplot[fill=cMlpSK]     coordinates {(C1,18) (C2,20) (C3,35)};
          \legend{ADALINE SB, ADALINE CB, Multicamadas SB, Multicamadas CB}
        \end{axis}
      \end{tikzpicture}

      \vspace{0.2em}
      {\tiny $\downarrow$ menos FN = melhor clinicamente}
    \end{column}
  \end{columns}
\end{frame}

% =======================================================================
% SLIDE 11 — RESULTADOS BANCO 2 (texto)
% =======================================================================
\begin{frame}{Resultados — Banco 2: Heart Failure}
  \centering\small
  \textit{299 pacientes: 96 óbitos | 203 sobreviventes}
  \vspace{0.3em}

  \renewcommand{\arraystretch}{1.35}
  \begin{tabular}{lllll}
    \toprule
    \textbf{Modelo} & \textbf{C1} & \textbf{C2} & \textbf{C3} & \textbf{FN mín.} \\
    \midrule
    Adaline Sem Biblioteca & 68,89\% / 17\,FN & 72,22\% / 17\,FN & \textbf{75,56\% / 16\,FN} & 16 (C3) \\
    Adaline Com Biblioteca (=) & 68,89\% / 17\,FN & 68,89\% / 17\,FN & 68,89\% / 17\,FN & 17 (=) \\
    Multicamadas Sem Biblioteca     & \cellcolor{ufvjmbluelight}\textbf{78,89\% / 9\,FN} & 75,56\% / 14\,FN & 71,11\% / 23\,FN & 9 (C1) \\
    Multicamadas Com Biblioteca     & \textbf{73,33\% / 18\,FN} & 68,89\% / 17\,FN & 63,33\% / 8\,FN$^*$ & 17 (C2) \\
    \bottomrule
  \end{tabular}

  {\tiny $^*$8\,FN com 63,3\% de acurácia: modelo classifica quase todos como doentes — clinicamente não confiável}

  \vfill
  \begin{columns}[T]
    \begin{column}{0.50\textwidth}
      \begin{block}{Maior acurácia}
        Multicamadas Sem Biblioteca C1: \textbf{78,89\%}
      \end{block}
    \end{column}
    \begin{column}{0.50\textwidth}
      \begin{alertblock}{Melhor clinicamente}
        Multicamadas Sem Biblioteca C1: \textbf{9 FN}
      \end{alertblock}
    \end{column}
  \end{columns}
  \vspace{0.2em}
  \centering{\small\textbf{Melhor modelo geral:} Multicamadas Sem Biblioteca C1 — melhor equilíbrio entre acurácia (78,89\%) e FN mínimo (9)}
\end{frame}

% =======================================================================
% SLIDE 12 — RESULTADOS BANCO 2 (gráficos)
% =======================================================================
\begin{frame}{Banco 2 — Visualização dos Resultados}
  \begin{columns}[c]
    % ---- gráfico acurácia ----
    \begin{column}{0.50\textwidth}
      \centering
      {\small\textbf{Acurácia por Configuração (\%)}}
      \vspace{0.15em}

      \begin{tikzpicture}
        \begin{axis}[
          ybar,
          width=1.05\columnwidth,
          height=5.6cm,
          symbolic x coords={C1,C2,C3},
          xtick=data,
          xticklabel style={font=\scriptsize},
          ymin=58, ymax=88,
          ytick={60,65,70,75,80,85},
          ylabel={\%},
          ylabel style={font=\scriptsize,rotate=-90,at={(ticklabel cs:0.5)},anchor=south},
          yticklabel style={font=\scriptsize},
          bar width=5.5pt,
          legend style={at={(0.5,-0.22)},anchor=north,legend columns=2,font=\tiny,
                        draw=none,fill=none,column sep=2pt},
          legend cell align=left,
          every axis plot/.append style={draw=none},
          axis line style={draw=ufvjmgray,thin},
          tick style={draw=ufvjmgray,thin},
          ymajorgrids=true,
          grid style={dotted,gray!40},
          clip=false,
          nodes near coords={\pgfmathprintnumber[fixed,precision=1]{\pgfplotspointmeta}},
          nodes near coords style={font=\tiny,inner sep=0pt,rotate=90,anchor=west,xshift=2pt},
        ]
          \addplot[fill=cAdalineNP] coordinates {(C1,68.89) (C2,72.22) (C3,75.56)};
          \addplot[fill=cAdalineSK] coordinates {(C1,68.89) (C2,68.89) (C3,68.89)};
          \addplot[fill=cMlpNP]     coordinates {(C1,78.89) (C2,75.56) (C3,71.11)};
          \addplot[fill=cMlpSK]     coordinates {(C1,73.33) (C2,68.89) (C3,63.33)};
          \legend{ADALINE SB, ADALINE CB, Multicamadas SB, Multicamadas CB}
        \end{axis}
      \end{tikzpicture}
    \end{column}
    % ---- gráfico FN ----
    \begin{column}{0.50\textwidth}
      \centering
      {\small\textbf{Falsos Negativos por Configuração}}
      \vspace{0.15em}

      \begin{tikzpicture}
        \begin{axis}[
          ybar,
          width=1.05\columnwidth,
          height=5.6cm,
          symbolic x coords={C1,C2,C3},
          xtick=data,
          xticklabel style={font=\scriptsize},
          ymin=0, ymax=28,
          ytick={0,5,10,15,20,25},
          ylabel={FN},
          ylabel style={font=\scriptsize,rotate=-90,at={(ticklabel cs:0.5)},anchor=south},
          yticklabel style={font=\scriptsize},
          bar width=5.5pt,
          legend style={at={(0.5,-0.22)},anchor=north,legend columns=2,font=\tiny,
                        draw=none,fill=none,column sep=2pt},
          legend cell align=left,
          every axis plot/.append style={draw=none},
          axis line style={draw=ufvjmgray,thin},
          tick style={draw=ufvjmgray,thin},
          ymajorgrids=true,
          grid style={dotted,gray!40},
          nodes near coords,
          nodes near coords style={font=\tiny,inner sep=1pt},
          nodes near coords align=above,
        ]
          \addplot[fill=cAdalineNP] coordinates {(C1,17) (C2,17) (C3,16)};
          \addplot[fill=cAdalineSK] coordinates {(C1,17) (C2,17) (C3,17)};
          \addplot[fill=cMlpNP]     coordinates {(C1,9)  (C2,14) (C3,23)};
          \addplot[fill=cMlpSK]     coordinates {(C1,18) (C2,17) (C3,8)};
          \legend{ADALINE SB, ADALINE CB, Multicamadas SB, Multicamadas CB}
        \end{axis}
      \end{tikzpicture}
    \end{column}
  \end{columns}
  \vspace{-0.3em}
  \centering{\tiny Overfitting visível: Multicamadas SB e Multicamadas CB pioram progressivamente com mais épocas}
\end{frame}

% =======================================================================
% SLIDE 13 — BINARIZADOS vs CONTÍNUOS
% =======================================================================
\section{Análise Complementar}

\begin{frame}{Experimento Complementar: Binarizados \textit{vs} Contínuos}
  \textbf{Hipótese:} A binarização favoreceria o ADALINE por simplificar o espaço de decisão.
  \vspace{0.3em}

  \begin{columns}[T]
    \begin{column}{0.44\textwidth}
      \begin{block}{Banco 1 — resultado próximo}
        \small Acurácias similares. A binarização não trouxe ganho significativo, mas também não prejudicou.
      \end{block}

      \vspace{0.3em}
      \begin{alertblock}{Banco 2 — resultado revelador}
        \small ADALINE SB subiu de 68,9\% para \textbf{82,2\%}\\
        FN: de 17 para \textbf{11}\\[0.15em]
        Multicamadas CB: caiu para \textbf{45,6\%} (instável)
      \end{alertblock}
    \end{column}
    \begin{column}{0.56\textwidth}
      \centering
      {\small\textbf{Banco 2 — Acurácia: Binarizados \textit{vs} Contínuos}}
      \vspace{0.15em}

      \begin{tikzpicture}
        \begin{axis}[
          ybar,
          width=\columnwidth,
          height=5.5cm,
          symbolic x coords={ADALINE SB, ADALINE CB, Multicamadas SB, Multicamadas CB},
          xtick=data,
          xticklabel style={font=\tiny,rotate=20,anchor=east},
          ymin=30, ymax=90,
          ytick={30,40,50,60,70,80,90},
          ylabel={\%},
          ylabel style={font=\scriptsize,rotate=-90,at={(ticklabel cs:0.5)},anchor=south},
          yticklabel style={font=\scriptsize},
          bar width=7pt,
          legend style={at={(0.5,-0.26)},anchor=north,legend columns=2,font=\tiny,
                        draw=none,fill=none},
          legend cell align=left,
          every axis plot/.append style={draw=none},
          axis line style={draw=ufvjmgray,thin},
          tick style={draw=ufvjmgray,thin},
          ymajorgrids=true,
          grid style={dotted,gray!40},
          clip=false,
          nodes near coords={\pgfmathprintnumber[fixed,precision=1]{\pgfplotspointmeta}},
          nodes near coords style={font=\tiny,inner sep=0pt,rotate=90,anchor=west,xshift=2pt},
        ]
          \addplot[fill=cBin]  coordinates {(ADALINE SB,68.89) (ADALINE CB,68.89) (Multicamadas SB,78.89) (Multicamadas CB,73.33)};
          \addplot[fill=cCont] coordinates {(ADALINE SB,82.2)  (ADALINE CB,82.2)  (Multicamadas SB,78.89) (Multicamadas CB,45.6)};
          \legend{Binarizados, Contínuos}
        \end{axis}
      \end{tikzpicture}
    \end{column}
  \end{columns}
  \vspace{-0.2em}
  \centering{\tiny A hipótese não se confirmou: o ADALINE se beneficiou \textbf{mais} dos dados contínuos neste contexto}
\end{frame}

% =======================================================================
% SLIDE 14 — COMPARAÇÃO MANUAL vs BIBLIOTECA
% =======================================================================
\begin{frame}{Validação: Implementação Manual \textit{vs} Biblioteca}
  \begin{columns}[T]
    \begin{column}{0.38\textwidth}
      \begin{itemize}\setlength{\itemsep}{6pt}
        \item \textbf{ADALINE SB} \textit{vs} \textbf{ADALINE CB}:\\
              Resultados praticamente idênticos
        \item \textbf{Multicamadas SB} \textit{vs} \textbf{Multicamadas CB}:\\
              Divergências maiores
      \end{itemize}
      \vspace{0.5em}
      \begin{block}{Destaque}
        \small Multicamadas SB \textbf{superou} o Multicamadas CB em C3 (87,68\% \textit{vs} 82,61\%) — implementação manual \textbf{competitiva}.
      \end{block}
    \end{column}
    \begin{column}{0.62\textwidth}
      \begin{columns}[c]
        \begin{column}{0.50\textwidth}
          \centering
          {\tiny\textbf{ADALINE — SB \textit{vs} CB}}\\[-0.1em]
          {\tiny\textit{\color{gray}Resultados praticamente idênticos}}
          \vspace{0.15em}

          \begin{tikzpicture}
            \begin{axis}[
              ybar,
              width=1.05\columnwidth,
              height=5cm,
              symbolic x coords={C1,C2,C3},
              xtick=data,
              xticklabel style={font=\tiny},
              ymin=84, ymax=89,
              ytick={84,85,86,87,88,89},
              ylabel={\%},
              ylabel style={font=\tiny,rotate=-90,at={(ticklabel cs:0.5)},anchor=south},
              yticklabel style={font=\tiny},
              bar width=6pt,
              legend style={at={(0.5,-0.18)},anchor=north,legend columns=2,font=\tiny,
                            draw=none,fill=none},
              legend cell align=left,
              every axis plot/.append style={draw=none},
              axis line style={draw=ufvjmgray,thin},
              tick style={draw=ufvjmgray,thin},
              ymajorgrids=true,
              grid style={dotted,gray!40},
              clip=false,
              nodes near coords={\pgfmathprintnumber[fixed,precision=1]{\pgfplotspointmeta}},
              nodes near coords style={font=\tiny,inner sep=0pt,rotate=90,anchor=west,xshift=2pt},
            ]
              \addplot[fill=cBin]  coordinates {(C1,85.87) (C2,86.23) (C3,86.96)};
              \addplot[fill=cCont] coordinates {(C1,85.87) (C2,86.23) (C3,86.23)};
              \legend{Sem Bib., Com Bib.}
            \end{axis}
          \end{tikzpicture}
        \end{column}
        \begin{column}{0.50\textwidth}
          \centering
          {\tiny\textbf{Multicamadas — SB \textit{vs} CB}}\\[-0.1em]
          {\tiny\textit{\color{gray}Divergência maior em C3}}
          \vspace{0.15em}

          \begin{tikzpicture}
            \begin{axis}[
              ybar,
              width=1.05\columnwidth,
              height=5cm,
              symbolic x coords={C1,C2,C3},
              xtick=data,
              xticklabel style={font=\tiny},
              ymin=80, ymax=92,
              ytick={80,82,84,86,88,90,92},
              ylabel={\%},
              ylabel style={font=\tiny,rotate=-90,at={(ticklabel cs:0.5)},anchor=south},
              yticklabel style={font=\tiny},
              bar width=6pt,
              legend style={at={(0.5,-0.18)},anchor=north,legend columns=2,font=\tiny,
                            draw=none,fill=none},
              legend cell align=left,
              every axis plot/.append style={draw=none},
              axis line style={draw=ufvjmgray,thin},
              tick style={draw=ufvjmgray,thin},
              ymajorgrids=true,
              grid style={dotted,gray!40},
              clip=false,
              nodes near coords={\pgfmathprintnumber[fixed,precision=1]{\pgfplotspointmeta}},
              nodes near coords style={font=\tiny,inner sep=0pt,rotate=90,anchor=west,xshift=2pt},
            ]
              \addplot[fill=cBin]  coordinates {(C1,82.97) (C2,86.23) (C3,87.68)};
              \addplot[fill=cCont] coordinates {(C1,87.32) (C2,88.04) (C3,82.61)};
              \legend{Sem Bib., Com Bib.}
            \end{axis}
          \end{tikzpicture}
        \end{column}
      \end{columns}
    \end{column}
  \end{columns}
\end{frame}

% =======================================================================
% SLIDE 15 — CONCLUSÃO
% =======================================================================
\section{Conclusão}

\begin{frame}{Conclusão}
  \begin{columns}[T]
    \begin{column}{0.50\textwidth}
      \textbf{Melhores resultados:}
      \begin{itemize}\setlength{\itemsep}{4pt}
        \item \textbf{Banco 1:} Multicamadas Sem Biblioteca C3 — \textbf{melhor modelo geral} (87,68\% / \textbf{15 FN})
        \item \textbf{Banco 2:} Multicamadas Sem Biblioteca C1 — melhor equilíbrio (\textbf{78,89\% / 9 FN})
      \end{itemize}
    \end{column}
    \begin{column}{0.50\textwidth}
      \textbf{Conclusões das comparações:}
      \begin{itemize}\setlength{\itemsep}{4pt}
        \item \textbf{Qual modelo?} O Multicamadas superou o ADALINE em ambos os bancos — a complexidade extra \textbf{trouxe ganho real}
        \item \textbf{Manual \textit{vs} Biblioteca?} Implementações manuais produziram resultados \textbf{competitivos} e, no Multicamadas C3, até \textbf{superiores}
        \item \textbf{Binarizar ou não?} A binarização \textbf{não trouxe vantagem}; os dados contínuos favoreceram o ADALINE no Banco 2
      \end{itemize}
    \end{column}
  \end{columns}
\end{frame}

% =======================================================================
% SLIDE 16 — OBRIGADO / PERGUNTAS
% =======================================================================
\begin{frame}[plain]
  \centering
  \vfill
  {\Huge\color{ufvjmblue}\textbf{Obrigado!}}
  \vfill
  {\large Lucas dos Santos Martins}\\[0.3em]
  {\normalsize Orientador: Prof.\ Dr.\ Leonardo Lana de Carvalho}\\[0.3em]
  {\small Universidade Federal dos Vales do Jequitinhonha e Mucuri}
  \vfill
  {\normalsize\itshape À disposição para perguntas.}
  \vfill
\end{frame}

\end{document}
